\subsection{\texorpdfstring{$\cBun_{r,d}^{\mathrm{ss}}$ 的上同调}{半稳定向量丛叠的上同调}}\label{src:9.4.6-subsection}
在划时代的论文 \cite{AB83} 中, Atiyah 与 Bott 将等变 Morse 理论应用于无穷维联络空间模无穷维规范变换群上的 Yang--Mills 泛函, 计算了上同调 $\mathrm H^*(\cBun_{r,d}^{\mathrm{ss}},\mathbb Q)$. 我们将以更代数化的方法概述他们的策略.

第一步是证明上同调环 $\mathrm H^*(\cBun_{r,d},\mathbb Q)$ 由重言类自由生成. 选取基 $1\in\mathrm H^0(C,\mathbb Q)$, $\gamma_1,\ldots,\gamma_{2g}\in\mathrm H^1(C,\mathbb Q)$, $\mathrm{pt}\in\mathrm H^2(C,\mathbb Q)$. 利用 K\"unneth 公式, 可将 $C\times\cBun_{r,d}$ 上万有向量丛 $\mathcal F_{\mathrm{univ}}$ 的 Chern 类写为
\[
\begin{aligned}
c_i(\mathcal F_{\mathrm{univ}})&=1\otimes a_i+\sum_{j=1}^{2g}\gamma_j\otimes b_{ij}+\mathrm{pt}\otimes f_i\\
&\in\mathrm H^*(C,\mathbb Q)\otimes\mathrm H^*(\cBun_{r,d},\mathbb Q)\cong\mathrm H^*(C\times\cBun_{r,d},\mathbb Q),
\end{aligned}
\]
其中 $a_i,b_{ij},f_i\in\mathrm H^*(\cBun_{r,d},\mathbb Q)$.
\begin{proposition}\label{src:9.4.46}
\[
\mathrm H^*(\cBun_{r,d},\mathbb Q)\cong\mathbb Q[a_1,\ldots,a_r]\otimes\bigwedge[b_{i,j}]_{\substack{i=1,\ldots,r\\j=1,\ldots,2g}}\otimes\mathbb Q[f_2,\ldots,f_r].
\]
\end{proposition}
\begin{proof}
见 \cite[\S1--2]{AB83} 和 \cite[定理 4.7]{Hei10}.
\end{proof}
由此可以计算 Poincar\'e 多项式 $p(\cBun_{r,d},t):=\sum_{d=0}^{\infty}\dim\mathrm H^d(\cBun_{r,d},\mathbb Q)t^d$.
\begin{corollary}\label{src:9.4.47}
\[
p(\cBun_{r,d},t)=\frac{\prod_{i=1}^{n}(1+t^{2i-1})^{2g}}{\prod_{i=1}^{n}(1-t^{2i})\prod_{i=1}^{n-1}(1-t^{2i})}.
\]
\end{corollary}
第二步是利用 $\cBun_{r,d}$ 的 Harder--Narasimhan 分层 (\bookref{9.3.15}) 与 Kirwan 满射性 (\bookref{8.7.13}), 推导半稳定轨迹 $\cBun_{r,d}^{\mathrm{ss}}$ 的 Poincar\'e 多项式的递推公式. 为建立记号, 回忆若 $0=F_0\subseteq F_1\subseteq\cdots\subseteq F_l=F$ 是向量丛 $F$ 的 Harder--Narasimhan 滤过 (\bookref{9.1.21}), 则 $F$ 的 Harder--Narasimhan 多边形 $\mathcal P$ 是点 $(-\deg F_i,\mathrm{rk}\,F_i)$ 的凸包. Harder--Narasimhan 分层 (\bookref{9.3.15}) 是分解 $\cBun_{r,d}=\coprod_{\mathcal P}\cBun_{r,d}^{\mathcal P}$, 其中 $\cBun_{r,d}^{\mathcal P}$ 是参数化 Harder--Narasimhan 多边形为 $\mathcal P$ 的向量丛的局部闭子叠. 在凸多边形的包含偏序下, $\cBun_{r,d}^{\leq\mathcal P}:=\bigcup_{\mathcal P'\subseteq\mathcal P}\cBun_{r,d}^{\mathcal P'}$ 是开子叠. Kirwan 满射性 (\bookref{8.7.13} 和 \bookref{8.7.15}) 表明 Thom--Gysin 序列
\[
\cdots\longrightarrow\mathrm H^{*-d_{\mathcal P}}(\cBun_{r,d}^{\mathcal P})\longrightarrow\mathrm H^*(\cBun_{r,d}^{\leq\mathcal P})\longrightarrow\mathrm H^*(\cBun_{r,d}^{<\mathcal P})\longrightarrow\cdots
\]
分裂. 所以 $\mathrm H^*(\cBun_{r,d}^{\mathrm{ss}})$ 是 $\mathrm H^*(\cBun_{r,d})$ 的商, 也由重言类生成. 令 $d_{\mathcal P}$ 为 $\cBun_{r,d}^{\mathcal P}$ 在 $\cBun_{r,d}$ 中的余维. Kirwan 满射性还给出 Poincar\'e 多项式的递推公式 (另见 \cite[定理 10.10]{AB83}):
\begin{proposition}\label{src:9.4.48}
\begin{equation}\label{src:9.4.49}
p(\cBun_{r,d}^{\mathrm{ss}},t)=p(\cBun_{r,d},t)-\sum_{\mathcal P}p(\cBun_{r,d}^{\mathcal P},t)t^{d_{\mathcal P}}.
\end{equation}
此外, 若 $(r_i',d_i')=(r_i-r_{i-1},d_i-d_{i-1})$ 表示 $F\in\cBun_{r,d}^{\mathcal P}$ 的 Harder--Narasimhan 滤过的半稳定因子 $F_i/F_{i-1}$ 的秩与次数, 则
\begin{equation}\label{src:9.4.50}
p(\cBun_{r,d}^{\mathcal P},t)=\prod_i p(\cBun_{r_i',d_i'}^{\mathrm{ss}},t).
\end{equation}
\end{proposition}
由递推公式~\bookeqref{9.4.49} 和~\bookeqref{9.4.50} 推导 $p(\cBun_{r,d}^{\mathrm{ss}},t)$ 的显式公式是一个困难的组合问题. 秩为 $2$ 时, Zagier 给出了优美的解答 (见 \cite{Zag96} 和 \cite[定理 5.16]{Hei10}):
\[
p(\cBun_{2,d}^{\mathrm{ss}},t)=\begin{cases}
\dfrac{(1+t)^{2g}}{(1-t^2)^2(1-t^4)}\big((1+t^3)^{2g}-t^{2g}(1+t)^{2g}\big),&d\text{ 为奇数},\\[6pt]
\dfrac{(1+t)^{2g}}{(1-t^2)^2(1-t^4)}\big((1+t^3)^{2g}-t^{2g+2}(1+t)^{2g}\big),&d\text{ 为偶数}.
\end{cases}
\]
最后, 若 $r$ 与 $d$ 互素, 则 $\cBun_{r,d}^{\mathrm{ss}}\to M_{r,d}^{\mathrm{ps}}$ 是平凡 $\mathbb G_m$-gerbe (推论~\bookref{9.4.35}). 因而 $\mathrm H^*(\cBun_{r,d}^{\mathrm{ss}},\mathbb Q)=\mathrm H^*(M_{r,d}^{\mathrm{ps}},\mathbb Q)[z]$, 可据此推出 $M_{r,d}^{\mathrm{ps}}$ 的 Poincar\'e 多项式 \cite[公式 (9.3)]{AB83}. 当秩为 $2$ 且 $d$ 为奇数时, 有
\[
p(M_{2,d}^{\mathrm{ps}},t)=\frac{(1+t)^{2g}}{(1-t^2)(1-t^4)}\big((1+t^3)^{2g}-t^{2g+2}(1+t)^{2g}\big).
\]
Atiyah 与 Bott 的计算此后被重新研究, 并推广到许多其他情形, 例如 \cite{Kir86,HS10,BDN+25}.
